% !TEX root = ../../main-es.tex
% appendix/es/C-reproducibilidad.tex
\chapter{Protocolo de reproducibilidad}
\label{app:reproducibilidad}

\porque{Rigor analítico (EICT 20\%).}{%
  Lo que un tercero necesita para recrear los resultados del cap. \ref{cap:experimentos}.
  Fallo típico: omisión de versiones o seeds.}

\section{Software}
\porque{Versiones pinneadas.}{\OASys{} commit hash; versión de Rust; versión de MLX; versiones de LangGraph para el brazo C.}
\porque{Entorno de análisis.}{%
  El protocolo de §\ref{sec:protocolo} está enunciado de forma independiente de la
  herramienta —cualquier entorno que ajuste un GLMM binomial con efectos aleatorios
  cruzados sirve—; aquí se pinnea la elección concreta: lenguaje, librerías de ajuste
  (frecuentista y MCMC), versiones exactas y semilla del generador pseudoaleatorio del
  análisis, distinta de las seeds de ejecución del cap. \ref{cap:experimentos}.
  Los scripts de análisis se versionan junto al arnés.}

\section{Hardware}
M4 Pro, 24\,GB. Sin cómputo externo (a confirmar con el director, ver Gate 2 paso 9).

\section{Comandos}
\porque{Marcador.}{Secuencia exacta para instanciar los $k$ seeds y los $n=1..12$ niveles por brazo.}

\section{Semillas}
$k$ seeds fijas, reportadas y compartidas por los tres brazos (diseño apareado:
misma seed y misma tarea en A, B y C).
\porque{$k$ no es un número arbitrario.}{%
  El valor de $k$ sale del cálculo de potencia por simulación de §\ref{sec:protocolo}
  sobre el GLMM primario, no de un mínimo elegido a ojo. Registrar aquí, antes de
  correr: $\alpha$, potencia objetivo
  $1-\beta$, el OR mínimo detectable en la interacción, y el $k$ resultante.
  Reemplazar el marcador «$k \geq 5$» del borrador por ese valor.}

\section{Timeout y censura}
\porque{Parte del protocolo, no del análisis post-hoc.}{%
  Umbral de timeout por corrida y política de censura declarados en §\ref{sec:protocolo}.
  Registrar el umbral exacto en segundos y la tasa de censura observada por brazo.}
